\section{Contoh Terhitung}
\label{sec:contoh-elgamal}

Bagian ini memakai dua himpunan parameter untuk memperlihatkan prosedur
ElGamal secara utuh. Contoh pertama memakai $p = 2357$ dan pesan satu blok.
Contoh kedua memakai $p = 2273$ dan pesan yang dipecah menjadi dua blok.

\begin{example}[Pembangkitan kunci, enkripsi, dan dekripsi dengan $p = 2357$]
\label{ex:elgamal-2357}
Bob membangkitkan kunci publik dan kunci privatnya. Alice akan mengenkripsi
pesan dengan menggunakan kunci publik Bob.

\textbf{(a) Pembangkitan kunci oleh Bob.} Misalkan $p = 2357$, $g = 2$, dan
$x = 1751$. Bilangan $g = 2$ adalah akar primitif dari $2357$, dan $x = 1751$
memenuhi syarat $2 \le x \le 2355$. Bob menghitung
% Koreksi: slide hlm. 14 mencetak modulus 2353 pada perhitungan ini; modulus
% yang benar adalah p = 2357. Nilai 2^1751 mod 2353 = 1008, bukan 1185.
\[ y = g^{x} \bmod p = 2^{1751} \bmod 2357 = 1185 . \]
Kunci publik Bob adalah $(y = 1185, g = 2, p = 2357)$, sedangkan kunci
privatnya $(x = 1751, p = 2357)$. Bob memberitahukan kunci publik itu kepada
Alice.

\textbf{(b) Enkripsi oleh Alice.} Misalkan pesan $m = 2035$. Nilai $m$ berada
di dalam selang $[0, 2356]$. Alice memilih bilangan acak $k = 1520$, lalu
menghitung
\begin{align*}
a &= g^{k} \bmod p = 2^{1520} \bmod 2357 = 1430, \\
b &= y^{k} \cdot m \bmod p = 1185^{1520} \cdot 2035 \bmod 2357 = 697 .
\end{align*}
Jadi, cipherteks yang dihasilkan adalah $(a, b) = (1430, 697)$. Alice
mengirim cipherteks itu kepada Bob.

\textbf{(c) Dekripsi oleh Bob.} Bob memakai kunci privatnya $(x = 1751,
p = 2357)$ untuk menghitung
\[
\left(a^{x}\right)^{-1} = a^{\,p-1-x} \bmod p
= 1430^{\,2357-1-1751} \bmod 2357
= 1430^{\,605} \bmod 2357 = 872 ,
\]
lalu
\[
m = b \cdot \left(a^{x}\right)^{-1} \bmod p
= 697 \cdot 872 \bmod 2357 = 2035 .
\]
Bob mendapatkan kembali plainteks $m = 2035$ yang dikirim oleh Alice.
\end{example}

\begin{example}[Pesan ``HALO'' dengan $p = 2273$]
\label{ex:elgamal-halo}
Alice membangkitkan pasangan kuncinya. Ia memilih bilangan prima $p = 2273$,
akar primitif dari $p$ yaitu $g = 3$, dan $x = 243$. Alice kemudian menghitung
\[ y = g^{x} \bmod p = 3^{243} \bmod 2273 = 461 . \]
Jadi, kunci privat Alice adalah $(x = 243, p = 2273)$ dan kunci publiknya
adalah $(y = 461, g = 3, p = 2273)$.

Bob ingin mengirim plainteks ``HALO'' kepada Alice. Misalkan $A = 00$,
$B = 01$, $\dots$, $Z = 25$. Dengan aturan itu, pesan $m$ dikodekan menjadi
bilangan \code{07001114}. Pengodean setiap huruf diperlihatkan pada
Tabel~\ref{tab:pengodean-halo}.

\begin{table}[htbp]
    \centering
    \begin{tabular}{@{}l c c c c@{}}
    \toprule
    Huruf & H & A & L & O \\
    \midrule
    Kode  & 07 & 00 & 11 & 14 \\
    \bottomrule
    \end{tabular}
    \caption{Pengodean kata ``HALO'' menjadi bilangan dengan aturan $A = 00$
    sampai $Z = 25$, sehingga pesannya menjadi $m = 07001114$}
    \label{tab:pengodean-halo}
\end{table}

Bob memecah $m$ menjadi blok yang lebih kecil, yaitu blok-blok sepanjang empat
angka. Hasilnya adalah $m_{1} = 0700$ dan $m_{2} = 1114$. Kedua nilai itu masih
terletak di dalam selang $[0, 2272]$ agar transformasi menjadi satu-ke-satu.

\textbf{(i) Enkripsi $m_{1} = 0700$.} Bob memilih bilangan acak $k = 1463$.
Dengan kunci publik Alice $(y = 461, g = 3, p = 2273)$, ia menghitung
\begin{align*}
a &= g^{k} \bmod p = 3^{1463} \bmod 2273 = 1439, \\
b &= y^{k} \cdot m_{1} \bmod p = 461^{1463} \cdot 700 \bmod 2273 = 74 .
\end{align*}
Jadi, cipherteks untuk $m_{1}$ adalah $c_{1} = (1439, 74)$. Angka $0700$ pada
perhitungan itu sama dengan $700$.

\textbf{(ii) Enkripsi $m_{2} = 1114$.} Bob memilih bilangan acak $k = 2001$.
Dengan kunci publik yang sama, ia menghitung
\begin{align*}
a &= g^{k} \bmod p = 3^{2001} \bmod 2273 = 1220, \\
b &= y^{k} \cdot m_{2} \bmod p = 461^{2001} \cdot 1114 \bmod 2273 = 1682 .
\end{align*}
Jadi, cipherteks untuk $m_{2}$ adalah $c_{2} = (1220, 1682)$. Bob mengirim
cipherteks $(1439, 74)$ dan $(1220, 1682)$ kepada Alice.

\textbf{(iii) Dekripsi oleh Alice.} Alice mendekripsi kedua cipherteks itu
dengan kunci privatnya $(x = 243, p = 2273)$.

Untuk $c_{1} = (1439, 74)$:
\[
\left(a^{x}\right)^{-1} = 1439^{\,2273-1-243} \bmod 2273
= 1439^{\,2029} \bmod 2273 = 1791 ,
\]
\[ m_{1} = 74 \cdot 1791 \bmod 2273 = 700 = 0700 . \]

Untuk $c_{2} = (1220, 1682)$:
\[
\left(a^{x}\right)^{-1} = 1220^{\,2029} \bmod 2273 = 1125 ,
\]
\[ m_{2} = 1682 \cdot 1125 \bmod 2273 = 1114 . \]

Plainteks yang didekripsi adalah $m_{1} m_{2} = 07001114$. Bila dikodekan
menjadi huruf, rangkaian itu berbunyi ``HALO'', sama dengan plainteks yang
dikirim oleh Bob.
\end{example}
